\subsubsection{具体分析}

问题三要求分析作物关系和经济变量相关性对方案的影响，属于受控模拟条件下的优化与机理比较问题。本问采用关系调整趋势MILP和相关蒙特卡洛评估。

首先按作物类别构造替代与互补边，接着设置弱、中、强三档关系强度，然后生成最终经济变量相关情景，最后比较收益、下尾风险和41维作物累计面积结构。具体流程如图\ref{fig:analysis3_flow}所示。

\begin{figure}[H]
  \begin{center}
    \begin{tikzpicture}[x=1cm, y=0.7cm]
      \node[box] (n11) at (0, 0) {构造关系};
      \node[box] (n12) at (5, 0) {设置强度};
      \node[box] (n13) at (10, 0) {调整销量};
      \node[box] (n23) at (10, -3) {三档求解};
      \node[box] (n22) at (5, -3) {生成相关情景};
      \node[box] (n21) at (0, -3) {核验相关性};
      \node[box] (n31) at (0, -6) {配对评价};
      \node[box] (n32) at (5, -6) {宏观归因};
      \node[box] (n33) at (10, -6) {判断备用};
      \draw[arrow] (n11) -- (n12); \draw[arrow] (n12) -- (n13);
      \draw[arrow] (n13) -- ++(0,-1.0) -| (n23); \draw[arrow] (n23) -- (n22);
      \draw[arrow] (n22) -- (n21); \draw[arrow] (n21) -- (n31);
      \draw[arrow] (n31) -- (n32); \draw[arrow] (n32) -- (n33);
    \end{tikzpicture}
    \caption{问题三分析流程图}
    \label{fig:analysis3_flow}
  \end{center}
\end{figure}

\subsubsection{模型准备}

附件只有2023年单年截面，不能估计真实的交叉价格弹性。本文只按类别与用途构造模拟关系，并把所有边标记为模拟假设。已有研究曾用横截面相关系数讨论作物关系\upcite{zhao2025relation}，但相关性不能推出因果性，因此正文只讨论在这些设定下方案如何变化。

问题三保持问题二的土地约束、管理配置、销售规则和评价函数不变，关闭关系结构的问题二正式方案作为基线。这样两者差异只来自新增关系设置。

关系强度设置为弱、中、强三档，系数分别为0.15、0.35和0.55。弱档用于观察较小关系扰动，强档用于检验收益波动的放大边界，中档作为正式代表方案。三档共享同一套地块适宜性、轮作和管理约束，并使用相同数量的测试情景，因此收益差异可以归因于关系响应和相关结构的改变，而不是模型规模或约束口径的变化。

替代边按作物类别在同类作物间构造（同类蔬菜、同类粮食等互相竞争），互补边在豆类与共享适宜地块季次的非豆类作物间构造（豆科固氮养地）\upcite{deaton1980,varacca2020}。由于缺少多年销量、价格和亩产量联动数据，这些边、方向、强度及相关矩阵均不能解释为真实市场弹性。本文将替代、互补和经济变量相关性明确标记为模拟假设，只检验这些设定在题面变化区间内对方案和收益分布的影响。相关性核验直接作用于最终导出的销量、亩产量和价格数组，宏观结构判断则采用全乡村41种作物的七年累计面积向量。
